\documentclass[10pt,a4paper]{amsart}
\usepackage{amsmath,amssymb,amsthm,graphicx}
\usepackage{verbatim}
\newenvironment{reference}{\comment}{\endcomment}
%\newenvironment{reference}{}{}
\newenvironment{slogan}{\comment}{\endcomment}
\newenvironment{history}{\comment}{\endcomment}

% For commutative diagrams we use Xy-pic
\usepackage[all]{xy}

% We use 2cell for 2-commutative diagrams.
\xyoption{2cell}
% Theorem environments.
%
\theoremstyle{plain}
\newtheorem{theorem}[subsection]{Theorem}
\newtheorem{proposition}[subsection]{Proposition}
\newtheorem{lemma}[subsection]{Lemma}

\theoremstyle{definition}
\newtheorem{definition}[subsection]{Definition}
\newtheorem{example}[subsection]{Example}
\newtheorem{exercise}[subsection]{Exercise}
\newtheorem{situation}[subsection]{Situation}

\theoremstyle{remark}
\newtheorem{remark}[subsection]{Remark}
\newtheorem{remarks}[subsection]{Remarks}

\numberwithin{equation}{subsection}

% Macros
%
\def\lim{\mathop{\mathrm{lim}}\nolimits}
\def\colim{\mathop{\mathrm{colim}}\nolimits}
\def\Spec{\mathop{\mathrm{Spec}}}
\def\Hom{\mathop{\mathrm{Hom}}\nolimits}
\def\Ext{\mathop{\mathrm{Ext}}\nolimits}
\def\SheafHom{\mathop{\mathcal{H}\!\mathit{om}}\nolimits}
\def\SheafExt{\mathop{\mathcal{E}\!\mathit{xt}}\nolimits}
\def\Sch{\mathit{Sch}}
\def\Mor{\mathop{\mathrm{Mor}}\nolimits}
\def\Ob{\mathop{\mathrm{Ob}}\nolimits}
\def\Sh{\mathop{\mathit{Sh}}\nolimits}
\def\NL{\mathop{N\!L}\nolimits}
\def\CH{\mathop{\mathrm{CH}}\nolimits}
\def\proetale{{pro\text{-}\acute{e}tale}}
\def\etale{{\acute{e}tale}}
\def\QCoh{\mathit{QCoh}}
\def\Ker{\mathop{\mathrm{Ker}}}
\def\Im{\mathop{\mathrm{Im}}}
\def\Coker{\mathop{\mathrm{Coker}}}
\def\Coim{\mathop{\mathrm{Coim}}}

% Boxtimes
%
\DeclareMathSymbol{\boxtimes}{\mathbin}{AMSa}{"02}

%
% Macros for moduli stacks/spaces
%
\def\QCohstack{\mathcal{QC}\!\mathit{oh}}
\def\Cohstack{\mathcal{C}\!\mathit{oh}}
\def\Spacesstack{\mathcal{S}\!\mathit{paces}}
\def\Quotfunctor{\mathrm{Quot}}
\def\Hilbfunctor{\mathrm{Hilb}}
\def\Curvesstack{\mathcal{C}\!\mathit{urves}}
\def\Polarizedstack{\mathcal{P}\!\mathit{olarized}}
\def\Complexesstack{\mathcal{C}\!\mathit{omplexes}}
% \Pic is the operator that assigns to X its picard group, usage \Pic(X)
% \Picardstack_{X/B} denotes the Picard stack of X over B
% \Picardfunctor_{X/B} denotes the Picard functor of X over B
\def\Pic{\mathop{\mathrm{Pic}}\nolimits}
\def\Picardstack{\mathcal{P}\!\mathit{ic}}
\def\Picardfunctor{\mathrm{Pic}}
\def\Deformationcategory{\mathcal{D}\!\mathit{ef}}
\usepackage{cleveref}
\crefname{originaldefinition}{original definition}{original definitions}
\Crefname{originaldefinition}{Original definition}{Original definitions}
\makeatletter
\newlabel{algebra-definition-locally-free}{{00NW}{source}{}{sourcepointer-1}{}}
\newlabel{algebra-definition-locally-free@cref}{{[originaldefinition][0][]00NW}{[1][0][]source}{}{}{}}
\makeatother
\crefname{originallemma}{original lemma}{original lemmas}
\Crefname{originallemma}{Original lemma}{Original lemmas}
\makeatletter
\newlabel{algebra-lemma-finite-projective}{{00NX}{source}{}{sourcepointer-2}{}}
\newlabel{algebra-lemma-finite-projective@cref}{{[originallemma][0][]00NX}{[1][0][]source}{}{}{}}
\makeatother
\crefname{originallemma}{original lemma}{original lemmas}
\Crefname{originallemma}{Original lemma}{Original lemmas}
\makeatletter
\newlabel{algebra-lemma-map-cannot-be-injective}{{05WI}{source}{}{sourcepointer-3}{}}
\newlabel{algebra-lemma-map-cannot-be-injective@cref}{{[originallemma][0][]05WI}{[1][0][]source}{}{}{}}
\makeatother
\title{A projective ideal that is not locally free}
\author{The Stacks Project Authors}
\date{stacks-project:examples.tex:local-snapshot-20261012; GNU FDL 1.2 or later}
\begin{document}
\maketitle
\noindent\textit{Editorial note.} All delivered mathematical text is the author's own, copied from the pinned Stacks Project chapter file examples.tex; no statement or proof step was rewritten, reordered or removed.\par\medskip
\noindent\textit{Editorial note.} The counted proof is the author's complete argument: The projectivity lemma's own proof (idempotent normalization, Hom colimit, Mittag-Leffler) is delivered in L48286 and is not repeated here.\par\medskip
\noindent\textit{Editorial note.} The label lemma-ideal-projective-not-locally-free names the counted statement, and the proof environment immediately following it in the source is delivered verbatim as its proof.\par\medskip
\noindent\textit{Editorial note.} The counted statement is carried by the author's own lemma environment in the chapter source.\par\medskip
\noindent\textit{Allowed standard prerequisites (editorial).}\begin{itemize}
\item Delivered card L48286 (same unit pool): an ideal generated by a countable collection of idempotents is projective as an R-module (author lemma-ideal-generated-by-idempotents-projective, source lines 2661-2683, with its complete Mittag-Leffler proof)
\item Context span 2685-2694: if P subset Q is an inclusion of R-modules with Q finite and P locally free, and Q is generated by N elements, then P is finite locally free with free parts of rank at most N (Algebra, map-cannot-be-injective), and P is a finite R-module (Algebra, finite-projective)
\item Context span 2696-2705: combining these, a non-finitely-generated ideal generated by a countable collection of idempotents is projective but not locally free; explicit example R = product over N of F\_2 with I generated by the idempotents e\_n = (1, ..., 1, 0, ...) with n ones
\end{itemize}\medskip
\ifdefined\hypertarget\hypertarget{sourcepointer-1}{}\fi
\noindent\textit{Source reference.} Original definition 00NW: sources/stacks-chow/tags.txt, source lines 874--874. Prerequisite pointer invoked at 2719; it names a standard fact of a sibling chapter that the counted argument uses, and every step of the counted proof built on it is retained verbatim in the delivered body.\par\medskip
\ifdefined\hypertarget\hypertarget{sourcepointer-2}{}\fi
\noindent\textit{Source reference.} Original lemma 00NX: sources/stacks-chow/tags.txt, source lines 875--875. Prerequisite pointer invoked at 2725; it names a standard fact of a sibling chapter that the counted argument uses, and every step of the counted proof built on it is retained verbatim in the delivered body.\par\medskip
\ifdefined\hypertarget\hypertarget{sourcepointer-3}{}\fi
\noindent\textit{Source reference.} Original lemma 05WI: sources/stacks-chow/tags.txt, source lines 7587--7587. Prerequisite pointer invoked at 2722; it names a standard fact of a sibling chapter that the counted argument uses, and every step of the counted proof built on it is retained verbatim in the delivered body.\par\medskip
\noindent
Suppose that $P \subset Q$ is an inclusion of $R$-modules with $Q$ a
finite $R$-module and $P$ locally free, see
Algebra, Definition \ref{algebra-definition-locally-free}.
Suppose that $Q$ can be generated by $N$ elements as an $R$-module.
Then it follows from
Algebra, Lemma \ref{algebra-lemma-map-cannot-be-injective}
that $P$ is finite locally free (with the free parts having rank
at most $N$). And in this case $P$ is a finite $R$-module, see
Algebra, Lemma \ref{algebra-lemma-finite-projective}.

\medskip\noindent
Combining this with the above we see that a non-finitely-generated
ideal which is generated by a countable collection of idempotents
is projective but not locally free. An explicit example is
$R = \prod_{n \in \mathbf{N}} \mathbf{F}_2$ and
$I$ the ideal generated by the idempotents
$$
e_n = (1, 1, \ldots, 1, 0, \ldots )
$$
where the sequence of $1$'s has length $n$.

\begin{lemma}
\label{lemma-ideal-projective-not-locally-free}
There exists a ring $R$ and an ideal $I$ such that $I$ is projective as
an $R$-module but not locally free as an $R$-module.
\end{lemma}

\begin{proof}
See above.
\end{proof}

\end{document}
